% What a Class Description Buys
%
% Written for the ACL style files (https://github.com/acl-org/acl-style-files),
% which are also the Overleaf "ACL Proceedings" template. Drop this file,
% references.bib and tables/ into that template and compile.
%
% Every figure in the prose comes from tables/macros.tex, generated by
% `make paper-tables` from results/. Do not type a number into the text: when
% the experiments are re-run the paper follows on its own, and a hardcoded
% number would silently become a lie.

\documentclass[11pt]{article}
\usepackage[review]{acl}          % drop the option for the camera-ready
\usepackage{times}
\usepackage{latexsym}
\usepackage{amsmath}          % \lVert, \tfrac, \big used in Section 3
\usepackage{booktabs}
\usepackage{graphicx}
\usepackage[T1]{fontenc}
\usepackage[utf8]{inputenc}
\usepackage{microtype}

% Generated by experiments/export_latex.py — do not edit by hand.
% Regenerate with: make paper-tables

\newcommand{\NEVAL}{299}
\newcommand{\OursTopOne}{36.1\%}
\newcommand{\OursTopThree}{67.9\%}
\newcommand{\OursFOne}{0.290}
\newcommand{\OursCI}{[30.8, 41.8]}
\newcommand{\TfidfTopOne}{26.1\%}
\newcommand{\TfidfFOne}{0.203}
\newcommand{\NErrors}{191}
\newcommand{\CorpusSize}{9{,}993}
\newcommand{\NSections}{21}


\title{What a Class Description Buys:\\
       Alignment, Not Wording, in Zero-Shot Taxonomy Classification}

% The [review] option to acl.sty anonymises this block automatically; it is
% filled in for the camera-ready.
\author{Dilay Dikbıyık \\
  Department of Computer Engineering, Kocaeli University \\
  \texttt{dilaydikbiyik@gmail.com} \\}

\begin{document}
\maketitle

% ─────────────────────────────────────────────────────────────────────────────
\begin{abstract}
Zero-shot text classification against a taxonomy is usually implemented by
embedding a short description of each class and ranking those descriptions
against the document. How much the wording of those descriptions matters is
rarely measured. We rewrite the class descriptions of a NACE Rev.~2 taxonomy
over German trade register texts, from category labels into definitions that
enumerate concrete activities, and gain \ContentEffect~points of top-1 accuracy.
A control shows the gain is not linguistic: rendering the same content in the
language of the documents is worth \LanguageEffect~points ($p =
\LanguageP$). We identify what does explain it — how far the rewrite moves each
class vector toward the centroid of the documents it must attract
($\rho = \AlignRho$, $p = \AlignP$ over \NCLASSES\ classes) — and show the
account is monotone across three corpora, one of which we predict before
measuring. The quantity cannot, however, be optimised per class: three selection
criteria, including direct optimisation of development accuracy, all fail to
generalise.
\end{abstract}

% ─────────────────────────────────────────────────────────────────────────────
\section{Introduction}

Assigning economic activity codes to free-text company descriptions is a
routine task for statistical offices, business registries and commercial data
providers, and a poorly supervised one: labelled examples are expensive,
taxonomies are revised, and the classes number in the tens or hundreds. The
standard remedy is zero-shot classification against the taxonomy itself. Each
class is represented by a short natural-language description, both description
and document are embedded by a sentence encoder, and the document is assigned to
the nearest class.

The class descriptions in such a pipeline are usually inherited rather than
designed. They come from the taxonomy's own published labels, from an internal
glossary, or from whoever set the system up. Their wording is treated as a
detail of configuration rather than as a modelling decision, and we are not
aware of a study that measures how much it costs to get wrong. This paper
argues it is the largest single lever in the pipeline, and identifies the
quantity that governs it.

Our contributions are:
\begin{itemize}
  \item Rewriting class descriptions from category labels into definitions that
        enumerate concrete activities gains \ContentEffect~points of top-1
        accuracy on NACE section assignment, validated on a held-out half of the
        evaluation set that was never inspected during the rewrite.
  \item A control separating content from language. Sixteen of twenty-one
        descriptions were in a language none of the documents were written in,
        which is an attractive explanation and the wrong one: the same content
        rendered in the documents' language is worth \LanguageEffect~points.
  \item A mechanism. The gain is predicted by how far the rewrite moves the
        class vector toward the centroid of its own documents
        ($\rho = \AlignRho$, $p = \AlignP$), the only candidate of five that
        holds within each corpus separately.
  \item A prospective confirmation on a third corpus, where the alignment change
        is computed and the prediction stated before accuracy is examined.
  \item A negative result with a measured cause: the quantity cannot be
        optimised per class, because the classifier's argmax compares across
        classes on a scale that depends on description style.
\end{itemize}

% ─────────────────────────────────────────────────────────────────────────────
\section{Related Work}

Zero-shot text classification by embedding label descriptions is a standard
construction, and sentence encoders trained with a contrastive sentence-level
objective \citep{reimers-gurevych-2019-sentence} — including multilingual ones
trained by distillation \citep{reimers-gurevych-2020-making} — are the usual
choice of encoder. Work in this line concentrates on the encoder, the scoring
function or the calibration of the output distribution. The label descriptions
themselves are typically taken as given.

The keyword-extraction literature is adjacent and instructive. KeyBERT
\citep{grootendorst-2020-keybert} scores candidate phrases by embedding
similarity to the document and admits a set of seed keywords intended to steer
extraction toward a domain; PatternRank
\citep{schopf-etal-2022-patternrank} constrains candidates by part of speech,
and PromptRank \citep{kong-etal-2023-promptrank} scores them with a prompted
language model. Purely statistical methods such as YAKE
\citep{campos-etal-2020-yake} remain competitive on short documents. The seed
mechanism is the closest analogue to what we study — a hand-written vocabulary
supplied to steer an embedding-based decision — and one of our findings is that
it is redundant once the class descriptions carry the same information.

% ─────────────────────────────────────────────────────────────────────────────
\section{Method}

A taxonomy assigns each of $K$ classes a name and, optionally, a description.
We form a class vector by encoding the concatenation of name and description,
and, in the configuration this repository shipped, averaging it with the
encoding of a hand-written seed keyword list:
\[
  v_c = \tfrac{1}{2}\big(f(\text{name}_c \Vert \text{desc}_c) + f(\text{seeds}_c)\big),
\]
where $f$ is the sentence encoder. A document $d$ is assigned to
$\arg\max_c \cos(f(d), v_c)$. Throughout the experiments that isolate the effect
of descriptions we drop the seed term, so that the quantity under test is the
only thing varying.

Our central measurement is the \emph{alignment} of a class vector with the
documents it should attract:
\[
  a(c) = \cos\!\big(v_c,\ \bar{d}_c\big),\qquad
  \bar{d}_c = \frac{\sum_{d \in D_c} f(d)}{\lVert \sum_{d \in D_c} f(d)\rVert},
\]
with $D_c$ the documents whose gold label is $c$. Rewriting a description
changes $v_c$ and therefore $a(c)$; we write $\Delta a(c)$ for that change.

% ─────────────────────────────────────────────────────────────────────────────
\section{Data}

\paragraph{NACE over trade register texts.}
\CorpusSize\ German company purpose statements from the Handelsregister, to be
assigned to one of the \NSections\ NACE Rev.~2 sections
\citep{eurostat-nace-2008}. The evaluation set is
\NEVAL\ documents sampled from that corpus, stratified by predicted section and
split between low-margin and representative cases; no gold label was used in the
selection. Labels were produced by a language model applying a written
annotation guideline and validated against a human pass: in a blind pilot a
second annotator agreed on 58\% of 50 documents ($\kappa = 0.542$); after the
disagreements were adjudicated into three explicit rules, agreement on a fresh
50-document sample was 80\% ($\kappa = 0.772$), rising to 100\% on the documents
the labeller had flagged as high-confidence. Verified answers were promoted over
the machine labels in the released set. We report this in full because it
bounds every number that follows.

An earlier version of this evaluation set consisted of \NEVAL\ hand-written
documents rather than corpus samples. It was measurably easier — median 116
characters against 175, and 3.3\% carrying the legal boilerplate that 24.5\% of
real entries carry — and the pipeline scored 80.0\% on it against 34.8\% on the
corpus-sampled set. We mention this because it is the reason the evaluation set
was rebuilt, and because it is a failure mode that is easy to reproduce
accidentally.

\paragraph{Two further corpora.}
20 Newsgroups (20 topical classes, English usenet posts) and Reuters-21578 (12
frequent categories, single-label documents only, English financial newswire).
The three differ in what their class labels look like: abstract economic
categories, readable topic names, and opaque codes such as \texttt{money-fx}
respectively.

\paragraph{Held-out protocol.}
The description rewrite was motivated by reading errors, which is a form of
fitting. It was therefore developed on a stratified development half of the
NACE evaluation set and reported on a held-out half that was not examined.

% ─────────────────────────────────────────────────────────────────────────────
\section{Experiments}

\subsection{Language or content?}

The rewritten descriptions changed two things at once: language, since sixteen
of twenty-one were originally in Turkish while every document is German, and
content, since terse category labels became definitions enumerating concrete
activities. Table~\ref{tab:descriptions} separates them. The control condition
renders the original labels in German by hand — same wording, same terseness, no
machine-translation noise, and the five sections already in German carried over
untouched.

% Generated by experiments/export_latex.py -- do not edit by hand.
% Regenerate with: make paper-tables

\begin{table}[t]
\centering
\small
\begin{tabular}{@{}>{\raggedright\arraybackslash}p{3.5cm}rrr@{}}
\toprule
Class descriptions & Top-1 & Top-3 & Words \\
\midrule
original (16 of 21 in Turkish, terse) & 28.4 & 62.5 & 10 \\
same content, rendered in German & 26.1 & 58.9 & 10 \\
\textbf{rewritten as NACE-style definitions} & 52.8 & 81.9 & 19 \\
\bottomrule
\end{tabular}
\caption{Class descriptions on \NEVAL\ NACE documents, with the sector vector built from the description alone. Rendering the same content in German is the control: it isolates language from content. Top-1 and Top-3 are percentages.}
\label{tab:descriptions}
\end{table}


Language alone is worth \LanguageEffect~points ($p = \LanguageP$). Content is
worth \ContentEffect. A crossed design varying document language and description
language independently agrees: the richer descriptions win in both document
languages, so the diagonal a matching effect would produce does not appear. The
effect also survives a change of encoder.

\subsection{What predicts which classes gain}

Table~\ref{tab:predictors} regresses the per-class gain on five candidate
predictors over \NCLASSES\ classes. The target is the share of \emph{available}
headroom a class captured rather than raw gain, since a class already near
ceiling cannot gain much and raw gain largely measures where a class started.

% Generated by experiments/export_latex.py — do not edit by hand.
% Regenerate with: make paper-tables

\begin{table}[t]
\centering
\small
\begin{tabular}{lrrrr}
\toprule
Predictor & $\rho$ & $p$ & NACE & 20NG \\
\midrule
\textbf{Movement toward own documents} & $+0.513$ & 0.003 & $+0.56$ & $+0.63$ \\
Alignment before rewriting & $-0.293$ & 0.104 & $-0.55$ & $-0.38$ \\
Confusability of the class vector & $-0.010$ & 0.957 & $+0.47$ & $-0.10$ \\
Lexical overlap of the class name & $+0.065$ & 0.726 & $+0.37$ & $-0.05$ \\
Words added & $-0.209$ & 0.251 & $-0.24$ & $+0.08$ \\
\bottomrule
\end{tabular}
\caption{Spearman correlation with the share of available headroom a class captured, over \NCLASSES\ classes. Only movement toward the class's own documents survives, and only it holds within each corpus separately.}
\label{tab:predictors}
\end{table}


Only $\Delta a(c)$ survives, and it is the only candidate that holds within each
corpus separately. Neither the lexical overlap between a class name and its
documents nor the confusability of the class vector predicts anything. Words
added correlates \emph{negatively}: the recipe is not to write more.

\subsection{Three corpora, one predicted in advance}

If $\Delta a$ is the mechanism it should also account for differences
\emph{between} corpora, and it should predict rather than explain.
Table~\ref{tab:corpora} reports all three. For Reuters, the alignment change was
computed and the prediction printed before accuracy was examined; the code that
enforces that ordering is released with the paper.

% Generated by experiments/export_latex.py -- do not edit by hand.
% Regenerate with: make paper-tables

\begin{table}[t]
\centering
\small
\setlength{\tabcolsep}{4pt}
\begin{tabular}{@{}llrr@{}}
\toprule
Corpus & Labels & $\Delta$align. & Gain \\
\midrule
NACE Rev. 2 & abstract categories & $+0.114$ & $+26.8$ \\
Reuters-21578 & opaque codes & $+0.028$ & $+4.8$ \\
20 Newsgroups & readable names & $-0.076$ & $+0.4$ \\
\bottomrule
\end{tabular}
\caption{The alignment change predicts the accuracy gained from elaborating class descriptions, monotonically across three corpora. Reuters was predicted before it was measured.}
\label{tab:corpora}
\end{table}


Elaboration is worth \ReutersGain~points on Reuters ($p < 0.001$) and
\NewsGain\ on 20 Newsgroups ($p = \NewsP$). The 20 Newsgroups definitions moved
their class vectors \emph{away} from their documents: twenty-three careful words
about baseball sit further from real usenet posts than the word ``Baseball''
does. Spelling a bare Reuters category code out into a readable name is worth a
further \ReutersCodeGain~points, the same identifier-to-name jump 20 Newsgroups
shows, and for the same reason.

\subsection{It cannot be optimised per class}

If $\Delta a$ were locally causal, choosing each class's description by it should
beat a fixed policy. It does not. Three criteria of increasing strength —
marginal alignment, a contrastive margin against the other classes, and greedy
search directly on development accuracy — all fail to generalise, the
contrastive one significantly so ($-3.6$ points, $p = 0.005$ on 20 Newsgroups).

Two causes, both measured. First, the styles do not share a similarity scale: in
a set where every other class carries an elaborated description, that half wins
69.8\% of argmax decisions on NACE and 35.2\% on 20 Newsgroups against a fair
share of 50\%, the bias running in the direction alignment predicts. Per-class
standardisation does not rescue it, costing 8.4 points on the best pure policy by
discarding class priors. Second, there are $2^K$ configurations for $K$ classes;
development accuracy rises under greedy search and held-out accuracy does not
follow.

The practical consequence is that description style should be uniform across
classes. A set of individually better descriptions can classify worse than a set
of consistently written ones.

\subsection{The seed keyword lists are redundant}

% Generated by experiments/export_latex.py — do not edit by hand.
% Regenerate with: make paper-tables

\begin{table}[t]
\centering
\small
\begin{tabular}{lrrr}
\toprule
Variant & Top-1 & $\Delta$ & F1-macro \\
\midrule
\textbf{Full system} & 80.0 & -- & 0.831 \\
\midrule
$-$ seeds in sector vector & 73.3 & $-6.7$ & 0.744 \\
$-$ description in sector vector & 76.7 & $-3.3$ & 0.770 \\
$-$ seed-guided extraction & 80.0 & $+0.0$ & 0.831 \\
$-$ six-stage keyword filter & 80.0 & $+0.0$ & 0.831 \\
+ cleaned text into the classifier & 86.7 & $+6.7$ & 0.833 \\
$\leftrightarrow$ mpnet-base-v2 encoder (768-dim) & 73.3 & $-6.7$ & 0.668 \\
$\leftrightarrow$ German translated to English first & 80.0 & $+0.0$ & 0.767 \\
\bottomrule
\end{tabular}
\caption{Component ablation. $\Delta$ is the change in Top-1 accuracy against the full system in percentage points.}
\label{tab:ablation}
\end{table}


With language-matched, informative descriptions in place, removing the
hand-written seed keyword lists entirely costs $+0.3$ points ($p = 1.000$). They
help by 4.6 points on the development half and hurt by 5.4 on the held-out half,
neither significant, so the defensible claim is that they add nothing reliable
rather than that they are harmful. Two further pipeline stages — seed-guided
keyword extraction and a six-stage keyword filter — show no effect in any
configuration we tested.

% Generated by experiments/export_latex.py — do not edit by hand.
% Regenerate with: make paper-tables

\begin{table*}[t]
\centering
\small
\begin{tabular}{lrrrrrr}
\toprule
System & Top-1 & 95\% CI & Top-3 & F1-macro & $\kappa$ & $p$ \\
\midrule
Random & 6.7 & [0.0, 16.7] & 26.7 & 0.033 & 0.009 & 0.000 \\
Majority class (oracle floor)$^{\dagger}$ & 13.3 & [3.3, 26.7] & 33.3 & 0.013 & 0.000 & 0.000 \\
TF-IDF → nearest NACE section & 76.7 & [60.0, 90.0] & 80.0 & 0.657 & 0.749 & 1.000 \\
Zero-shot embeddings (no taxonomy guidance) & 73.3 & [56.7, 86.7] & 86.7 & 0.744 & 0.714 & 0.500 \\
Unguided KeyBERT & 73.3 & [56.7, 86.7] & 86.7 & 0.744 & 0.714 & 0.500 \\
\textbf{Ours: taxonomy-guided} & 80.0 & [63.3, 93.3] & 96.7 & 0.831 & 0.786 & -- \\
\bottomrule
\end{tabular}
\caption{Section classification on the held-out evaluation set ($n{=}\NEVAL$). Intervals are percentile bootstrap over $10{,}000$ resamples; $p$ is an exact McNemar test against our full system on the same documents. $\dagger$ uses the gold label distribution and is an oracle floor, not a competing method.}
\label{tab:baselines}
\end{table*}


Table~\ref{tab:baselines} places the system against lexical and embedding
baselines. Top-3 accuracy of \OursTopThree\ is the operating point we would
argue for: the task admits genuine ambiguity, and a tool that proposes three
codes to a human coder is the realistic deployment.

% ─────────────────────────────────────────────────────────────────────────────
\section{Error Analysis}

We hand-coded fifty errors from the development half against a written codebook.

% Generated by experiments/export_latex.py — do not edit by hand.
% Regenerate with: make paper-tables

\begin{table}[t]
\centering
\small
\begin{tabular}{lr}
\toprule
Property of the error & Count \\
\midrule
recoverable in top3 & 5 \\
low margin & 4 \\
crowded top3 & 2 \\
outside top3 & 1 \\
\bottomrule
\end{tabular}
\caption{Automatically determined properties of the 6 misclassified documents. Manual categories are assigned against the codebook in the accompanying repository.}
\label{tab:errors}
\end{table}


Seed leakage is the largest category at 34\%, and the failures are lexical
rather than semantic: \emph{Bodenbelagsarbeiten} (floor covering, Construction)
is drawn to Mining because \emph{Boden} appears in that section's vocabulary;
\emph{Gebäudereinigung} leaves Administrative services for Construction on
\emph{Gebäude}. This converges with the ablation from the other direction and is
the second reason to drop the seed lists.

Ambiguous section definitions and genuinely multi-sector companies together
account for 48\% of errors and are not addressable by any change to this system.
NACE admits two defensible answers for a company that both manufactures and
installs windows. We read this as a rough ceiling on single-label accuracy for
the task, and as the strongest argument for reporting top-3.

% ─────────────────────────────────────────────────────────────────────────────
\section{Conclusion}

What a class description buys a zero-shot classifier is alignment with the
documents it must attract, and the wording of that description is a modelling
decision rather than a configuration detail. On a taxonomy whose class names
never appear in its documents, writing that vocabulary down is worth
\ContentEffect~points; on one whose class names are already the documents' own
words, the same work is worth nothing. The alignment change predicts both, and
predicted a third corpus before it was measured. It cannot, however, be turned
into a per-class selection rule, and we report the three ways that failed.

% ─────────────────────────────────────────────────────────────────────────────
\section*{Limitations}

The evaluation labels are model-assisted rather than gold: produced by a
language model applying a written guideline, validated against a human pass on
50 documents at $\kappa = 0.772$, with 50 of \NEVAL\ carrying human-verified
labels. Any figure in this paper should be read against that agreement.

Agreement was measured between a human and the labeller, not between two humans.
The inter-annotator figure reviewers usually ask for would require a second
annotator we did not have.

The alignment account rests on three corpora. Their alignment changes are well
separated ($+0.114$, $+0.028$, $-0.076$), which is what makes the ordering
visible, and equally what makes it a weak test of the account's behaviour near
zero. The class-level correlation, while significant and replicated within each
corpus, is $\rho = \AlignRho$: it orders classes, it does not determine them.

The class vectors, the documents and the alignment measure all come from one
encoder family. The description effect replicates across two encoders of
different size, but both are multilingual sentence-transformers trained by the
same distillation procedure.

We do not compare against an instruction-tuned language model asked to name a
class directly, which in 2026 is the obvious alternative implementation of
zero-shot taxonomy classification. The adapter is implemented in the released
code but was not run.

Finally, the per-class descriptions used in the elaborated conditions were
written by one author, who had seen the corpora. Style is therefore controlled
across conditions but not author bias.

\section*{Ethics Statement}

The trade register corpus consists of public records describing companies rather
than individuals, though registered company names sometimes contain the names of
their owners. The collected sample is not redistributed, because its
redistribution licence has not been established; the vocabulary and inverse
document frequency statistics derived from it are released instead, which are
aggregate counts and cannot be inverted to recover the records. The evaluation
set, the taxonomy, the annotation guideline and all experiment code are
released.

% Acknowledgements must stay empty in the review version: they are one of the
% commonest ways an anonymous submission is de-anonymised. Fill this in for the
% camera-ready only.
% \section*{Acknowledgements}

\bibliography{references}
\bibliographystyle{acl_natbib}

\end{document}
